% !TeX root = main.tex
% ================= DOCUMENT CLASS =================
% Hỗ trợ kích thước font 13pt
\documentclass[13pt,a4paper,oneside]{extreport}

% ================= GÓI CƠ BẢN & TIẾNG VIỆT =================
\usepackage[utf8]{inputenc}
\usepackage[utf8]{vietnam}
\usepackage{mathptmx} % Font Times New Roman
\usepackage{geometry}
\geometry{
    a4paper,
    left=30mm,
    right=20mm,
    top=25mm,
    bottom=25mm
}

% ================= GIÃN DÒNG & ĐOẠN VĂN =================
\usepackage{setspace}
\setstretch{1.3} % Tương đương giãn dòng 1.5 trong Word
\setlength{\parindent}{0cm} % Không thụt đầu dòng
\setlength{\parskip}{6pt}   % Khoảng cách giữa các đoạn văn

% ================= HÌNH ẢNH & ĐỒ HỌA =================
\usepackage{graphicx}
\graphicspath{{figures/}} % Tự động tìm trong thư mục figures/
\usepackage{tikz}
\usetikzlibrary{calc}

% ================= BẢNG BIỂU & TOÁN HỌC =================
\usepackage{array}
\usepackage{booktabs}
\usepackage{tabularx}
\usepackage{multirow}
\usepackage{amsmath,amssymb}
\usepackage{float}
\usepackage{url}
\usepackage{cite} % Định dạng trích dẫn chuẩn IEEE [1], [2]

% ================= CHÈN SOURCE CODE =================
\usepackage{listings}
\usepackage{xcolor}

\lstset{
    basicstyle=\ttfamily\fontsize{10}{12}\selectfont,
    backgroundcolor=\color{gray!10},
    frame=single,
    breaklines=true,
    keywordstyle=\bfseries\color{blue},
    commentstyle=\itshape\color{gray},
    stringstyle=\color{red!70!black},
    numbers=left,
    numberstyle=\tiny\color{gray},
    tabsize=4,
    showstringspaces=false
}

% ================= ĐỊNH DẠNG MỤC LỤC (TOCLOFT) =================
\usepackage[titles]{tocloft}
\setlength{\cftbeforechapskip}{6pt}   % Khoảng cách trước chương trong mục lục
\setlength{\cftbeforesecskip}{2pt}    % Khoảng cách giữa các mục section
\setlength{\cftbeforesubsecskip}{1pt} % Khoảng cách giữa các mục subsection
\setlength{\cftbeforetabskip}{3pt}    % Khoảng cách trong Danh mục bảng
\setlength{\cftbeforefigskip}{3pt}    % Khoảng cách trong Danh mục hình
\renewcommand{\cftchapleader}{\cftdotfill{\cftdotsep}} % Dấu chấm dẫn đến số trang

% ================= ĐÁNH SỐ TRANG (FANCYHDR) =================
\usepackage{fancyhdr}
\pagestyle{empty}
\fancypagestyle{plain}{
    \fancyhf{}
    \fancyfoot[C]{\fontsize{10}{12}\selectfont\thepage}
    \renewcommand{\headrulewidth}{0pt}
    \renewcommand{\footrulewidth}{0pt}
}

% ================= TIÊU ĐỀ TIẾNG VIỆT CHO DANH MỤC =================
\renewcommand{\contentsname}{MỤC LỤC}
\renewcommand{\listtablename}{DANH MỤC BẢNG BIỂU}
\renewcommand{\listfigurename}{DANH MỤC HÌNH ẢNH}
\renewcommand{\bibname}{TÀI LIỆU THAM KHẢO}
\renewcommand{\appendixname}{PHỤ LỤC}
\renewcommand{\tablename}{Bảng}
\renewcommand{\figurename}{Hình}

% ================= ĐỊNH DẠNG HEADING (TITLESEC) =================
\usepackage{titlesec}
\setcounter{tocdepth}{2}
\setcounter{secnumdepth}{3}

% Cấp 1: CHƯƠNG X. TÊN CHƯƠNG (14pt, In hoa, Đậm, Giữa)
\titleformat{\chapter}[block]
  {\normalfont\fontsize{14}{17}\bfseries\centering}
  {\MakeUppercase{CHƯƠNG\ \thechapter.}}{0.5em}{\MakeUppercase}
\titlespacing*{\chapter}{0pt}{0pt}{18pt}

% Cấp 2: 1.1. Tiêu đề mục (13pt, Đậm, Canh trái)
\titleformat{\section}
  {\normalfont\fontsize{13}{16}\bfseries}
  {\thesection.}{0.5em}{}
\titlespacing*{\section}{0pt}{12pt}{6pt}

% Cấp 3: 1.1.1. Tiêu đề tiểu mục (13pt, Nghiêng đậm, Canh trái)
\titleformat{\subsection}
  {\normalfont\fontsize{13}{16}\itshape\bfseries}
  {\thesubsection.}{0.5em}{}
\titlespacing*{\subsection}{0pt}{10pt}{4pt}

% Cấp 4: a. Tiêu đề mục con (13pt, Đậm, Canh trái)
\titleformat{\subsubsection}
  {\normalfont\fontsize{13}{16}\bfseries}
  {\alph{subsubsection}.}{0.5em}{}
\titlespacing*{\subsubsection}{0pt}{8pt}{2pt}

% ================= ĐỊNH DẠNG CHÚ THÍCH (CAPTION) =================
\usepackage{caption}
\captionsetup[table]{
    position=top,
    justification=raggedright,
    singlelinecheck=false,
    font={bf,small},
    skip=6pt,
    labelsep=period
}
\captionsetup[figure]{
    position=bottom,
    justification=centering,
    singlelinecheck=false,
    font={bf,small},
    skip=6pt,
    labelsep=period
}

% Lệnh chú thích nguồn tài liệu (12pt, Nghiêng, Canh phải)
\newcommand{\source}[1]{%
    \par\vspace{4pt}%
    {\raggedleft\fontsize{12}{14}\selectfont\itshape Nguồn: #1\par}%
    \vspace{6pt}%
}

% Vẽ khung viền đôi cho trang bìa
\newcommand{\drawborder}{
    \begin{tikzpicture}[remember picture, overlay]
        \draw[line width=1.8pt] 
            ($(current page.north west) + (20mm,-15mm)$) 
            rectangle 
            ($(current page.south east) + (-15mm,15mm)$);
        \draw[line width=0.6pt] 
            ($(current page.north west) + (21.5mm,-16.5mm)$) 
            rectangle 
            ($(current page.south east) + (-13.5mm,13.5mm)$);
    \end{tikzpicture}
}

% ================= THÔNG TIN DỰ ÁN / BÁO CÁO =================
\newcommand{\UnivName}{TRƯỜNG ĐẠI HỌC GIAO THÔNG VẬN TẢI TP. HỒ CHÍ MINH}
\newcommand{\InstituteName}{VIỆN CÔNG NGHỆ THÔNG TIN VÀ ĐIỆN, ĐIỆN TỬ}
\newcommand{\ReportType}{BÁO CÁO ĐỒ ÁN}
\newcommand{\ProjectTitle}{XÂY DỰNG VÀ PHÁT TRIỂN ỨNG DỤNG PETCAREAPP}
\newcommand{\MajorName}{CÔNG NGHỆ THÔNG TIN}
\newcommand{\SpecializationName}{CÔNG NGHỆ THÔNG TIN}
\newcommand{\CourseName}{Lập trình thiết bị di động}
\newcommand{\InstructorName}{ThS. Trương Quang Tuấn}
\newcommand{\StudentOne}{Dương Thành Đạt}
\newcommand{\StudentTwo}{Châu Tuấn Nguyên}
\newcommand{\StudentThree}{Châu Hoàng Quang Huy}
\newcommand{\LocationDate}{TP. Hồ Chí Minh, tháng 10 năm 2026}
